% Abstract. Every number carries a ledger id: TE figures E01a/E02a/E03a, the
% difference E04, the estimator spread E10, the BAB range E18, active sector
% E19a/E19b. Precision follows the ledger and must not be adjusted here.
\begin{abstract}
A published index methodology is not a complete specification. We rebuild two
low-volatility equity indices from free public data, one defined by a
constrained optimiser over a commercial factor risk model and one by a rank and
inverse-volatility rule, and ask how much of each is recoverable. The
transparent rule is almost fully recoverable: it tracks its fund at 1.62\%
annualised tracking error with a correlation of 0.9939. The optimiser is not,
at 5.83\%. The residual is not covariance estimation. Four estimators spanning
a sample matrix, a two-half-life structural model, constant-correlation
shrinkage and a Bai--Ng-selected principal-component model produce realised
volatilities 25 basis points apart under an identical constraint set, which is
what prior work reports for constrained portfolios generally. Resolving one
ambiguity in the published constraint document moves tracking error by 2.10
percentage points. The methodology exempts a single proprietary volatility
factor from its style bands; a replicator with a different factor set must
decide whether that exemption covers residual volatility alone or beta as well,
and exempting both brings tracking error to 3.73\%, replica volatility to
13.97\% against the fund's 13.45\%, and the beta to the fund to 1.001. On this
evidence the specification choice is worth roughly an order of magnitude more
than the estimator choice, for this index pair over this sample; we do not
establish that ordering in general. What the constraints deliver is control
rather than lower realised risk, holding the largest active industry weight
near its five per cent band where the heuristic averages 40.7\%. No portfolio
shows a six-factor alpha significant at the five per cent level, and
betting-against-beta loadings between 8.7 and 14.0 absorb what intercept there
is.
\end{abstract}
\noindent\textit{Keywords:} minimum volatility; index replication; factor risk model; covariance estimation; portfolio constraints.\par
\noindent\textit{JEL classification:} C13, C58, G11, G12, G17.
\vspace{8pt}
